\documentclass[11pt]{article}
\usepackage[margin=1in]{geometry}
\usepackage{amsmath,amssymb,amsthm}
\usepackage[hidelinks]{hyperref}
\newtheorem{theorem}{Theorem}
\newtheorem{lemma}[theorem]{Lemma}
\newtheorem{corollary}[theorem]{Corollary}
\newtheorem{proposition}[theorem]{Proposition}
\newcommand{\Q}{\mathbb Q}
\newcommand{\Fp}{\mathbb F_p}
\title{Residual lifting and parity for binomial factorial moments}
\author{}
\date{}
\begin{document}
\maketitle
For integers $a\ge0,d\ge1,n\ge1$ put
\[
 P_{n,a,d}(t)=\sum_{k=0}^n\binom nk(an+1)_{dk}t^k .
\]
These are the finite ordinary factorial-moment polynomials.
We identify their no-carry Newton residual with a smaller-index member
of the same family and obtain restrictions on rational factor degrees.
Real-root parity strengthens those restrictions. The results give
unbounded families of consecutive coprimality, including several quintic
families, but do not establish irreducibility or coprimality at every
order. No novelty or external-review claim is made.

\section{Real roots and parity}

\begin{lemma}\label{lem:real-root-parity}
For integers $a\ge0,d\ge1,n\ge1$ and every real $\beta>0$, put
\[
 P_{n,a,d,\beta}(t)=
 \sum_{k=0}^n\binom nk(an+\beta)_{dk}t^k.
\]
If $n$ is even, this polynomial is strictly positive on $\mathbb R$
and has no real roots. If $n$ is odd, its derivative is strictly
positive on $\mathbb R$, and it has exactly one real root, which is
simple. Every real-coefficient common divisor of
$P_{n,a,d,\beta}$ and $P_{n+1,a,d,\beta}$ therefore has even degree.
In particular, every rational factor of an even-index ordinary
polynomial $P_{n,a,d}$ has even degree.
\end{lemma}
\begin{proof}
For real $t$ the Gamma integral gives
\[
 P_{n,a,d,\beta}(t)=\frac1{\Gamma(an+\beta)}
 \int_0^\infty x^{an+\beta-1}(1+tx^d)^n e^{-x}\,dx .
 \tag{P1}
\]
This follows by expanding the finite power and integrating each monomial.
If $n$ is even, the integrand is nonnegative and positive on a nonempty
interval, proving strict positivity. If $n$ is odd, differentiating the
finite polynomial identity gives
\[
 P_{n,a,d,\beta}'(t)=\frac n{\Gamma(an+\beta)}
 \int_0^\infty x^{an+\beta+d-1}(1+tx^d)^{n-1}e^{-x}\,dx>0.
 \tag{P2}
\]
Its positive leading coefficient and odd degree then give exactly one
real root, and the derivative inequality makes it simple.
A common divisor cannot have any real root, because one member of a
consecutive pair has even index. A real polynomial of odd degree has
a real root, so such a divisor has even degree. The last assertion
uses the same argument for a factor of the positive even-index member.
\end{proof}

\section{Residual-index lifting under the no-carry hypothesis}

\begin{theorem}\label{thm:residual-index}
Let $p$ be prime, $r\ge1$, $b\ge1$, and suppose
\[
 (ab\bmod p)+db<p.
 \tag{1}
\]
Put
\[
 P=p^r,\quad n=bP,\quad
 s=d(1+p+\cdots+p^{r-1}),\quad \epsilon=(-1)^s\in\Fp^\times .
\]
The Newton residual of $P_{n,a,d}$ in the variable $Y=p^s t^P$ is exactly
\[
 {\cal R}(Y)=P_{b,a,d}(\epsilon Y)\pmod p.
 \tag{2}
\]
More precisely, every root has valuation $-s/P$, and every nonzero
factor $g(t)\in\Q_p[t]$ normalized by $g(0)=1$ has degree $P\ell$.
Writing $g=\sum a_kt^k$, its residual is
\[
 {\cal R}(g)(Y)=
 \sum_{j=0}^{\ell}\overline{a_{jP}/p^{sj}}Y^j,
 \tag{3}
\]
has exact degree $\ell$, and divides the residual in (2).
Consequently $\deg(g)/P$ is a sum of a submultiset of the irreducible
factor degrees of $P_{b,a,d}\bmod p$, with their multiplicities.
If that residual is irreducible of degree $b$, then
$P_{bp^r,a,d}$ is irreducible over $\Q_p$ and over $\Q$.
\end{theorem}

\begin{proof}
Condition (1) implies $db<p$, hence $b<p,d<p$.
The interval $(ab,ab+db]$ contains no multiple of $p$ even when $ab\ge p$.
For $0\le k\le n$, writing $c_k=\binom nk(an+1)_{dk}$, Legendre's formula
therefore gives
\[
 v_p((an+1)_{dk})=\sum_{i=1}^r\left\lfloor\frac{dk}{p^i}\right\rfloor.
 \tag{4}
\]
There are no contributions from higher powers, since the corresponding
full interval $(an,(a+d)n]$ contains no multiple of $p^{r+1}$.
For $0<k<n$, write $u=v_p(k)\le r$. One has
$v_p(n-j)=v_p(j)$ for $1\le j<n$:
the case $v_p(j)<r$ is immediate, and otherwise $j=hP$ with
$1\le h<b<p$. Thus $\binom{n-1}{k-1}$ is a unit, and
\[
 v_p\binom nk=r-u.
\]
If $P\nmid k$, the exact floors through $i=u$, and the strict estimates
$\lfloor x\rfloor>x-1$ thereafter, give
\[
 v_p(c_k)>k d\sum_{i=1}^rp^{-i}=ks/P.
\]
At $k=jP$, every floor is exact and $v_p(c_{jP})=sj$.
Thus all points lie on or above the endpoint chord, with equality only
at multiples of $P$. Also $s\not\equiv0\pmod p$.
At a hypothetical root of valuation greater than $-s/P$, the constant
term would be uniquely least; at a smaller valuation the leading term
would be uniquely least. This proves the root valuation assertion.

Let $U_p(m!)=m!/p^{v_p(m!)}$. Splitting the units in factorials into
blocks of length $p$, and pairing inverse elements of $\Fp^\times$, gives
\[
 U_p(N!)\equiv
 (-1)^{\lfloor N/p\rfloor}(N\bmod p)!\,
 U_p(\lfloor N/p\rfloor!)\pmod p.
\]
Iteration through the first $r$ zero digits gives, for every $m\ge0$,
\[
 U_p((mP)!)\equiv
 (-1)^{m(1+p+\cdots+p^{r-1})}U_p(m!)\pmod p.
 \tag{5}
\]
The right side contains the unit part of $m!$, not $m!$ itself.
In the ratio with $m=ab+dj$ and $m=ab$, condition (1) ensures
the intervening product has no factor divisible by $p$. Hence (5) gives
\[
 \frac{(an+1)_{djP}}{p^{sj}}
 \equiv\epsilon^j(ab+1)_{dj}\pmod p.
\]
The Frobenius identity
$(1+z)^{bP}\equiv(1+z^P)^b\pmod p$ gives
$\binom{bP}{jP}\equiv\binom bj$.
This proves (2), including the sign and the case $ab\ge p$.

For the factor assertion, write
$g(t)=\prod_{\nu=1}^h(1-t/\alpha_\nu)$ over a splitting field.
Every reciprocal root has valuation $s/P$, so
\[
 v_p(a_k)\ge ks/P,\qquad v_p(a_h)=hs/P.
 \tag{6}
\]
The last valuation is an integer because $a_h\in\Q_p^\times$.
Since $\gcd(s,P)=1$, $h=P\ell$, and (3) is an integral residual
of exact degree $\ell$ and constant term one.
If $g_1,g_2$ satisfy (6), then
\[
 {\cal R}(g_1g_2)={\cal R}(g_1){\cal R}(g_2).
 \tag{7}
\]
Indeed, at a coefficient of index $jP$, a cross term with first index
$k$ not divisible by $P$ has valuation at least
\[
 \left\lceil ks/P\right\rceil+
 \left\lceil(jP-k)s/P\right\rceil=sj+1.
\]
Only multiples of $P$ survive, and those give the usual product
convolution. Apply (7) to a factor and its complementary factor,
each normalized to have constant one. This proves the divisor and
submultiset assertions. An irreducible residual has no proper
nonconstant residual divisor, so it permits no proper factor.
\end{proof}

\begin{corollary}\label{cor:parity-residual}
Assume the hypotheses of Theorem~\ref{thm:residual-index}, and in addition
that $p$ is odd and $b$ is even. If the residual has no proper nonzero
even submultiset-sum of irreducible factor degrees, then
$P_{bp^r,a,d}$ is irreducible over $\Q$.

Equivalently, this condition on the residual factorization holds when
the residual is irreducible of degree $b$, or when it has exactly two
irreducible factors of odd degree, counted with multiplicity.
In particular it always holds for $b=2$.
\end{corollary}
\begin{proof}
Here $n=bp^r$ is even. Lemma~\ref{lem:real-root-parity} shows that
every nonconstant rational factor has even degree.
Since $p^r$ is odd, its quotient degree $\ell$ in Theorem~\ref{thm:residual-index}
must be even. The stated submultiset condition excludes every
proper factor.

For completeness, if a residual factorization of even total degree
has a proper irreducible factor of even degree, it already gives a
forbidden even submultiset. Otherwise all proper irreducible factors
have odd degree, so their number is even. If there are at least four,
two of them give a proper even sum. Exactly two odd factors do not
give a proper nonzero even sum. This proves the equivalence, with
multiplicities included.
\end{proof}

\begin{corollary}\label{cor:residual-adjacent-pairs}
Whenever Theorem~\ref{thm:residual-index} or
Corollary~\ref{cor:parity-residual} proves that $P_{n,a,d}$ is
irreducible over $\mathbb Q$, it is coprime to both
$P_{n-1,a,d}$ and $P_{n+1,a,d}$, with the first assertion
interpreted as coprimality with $P_0=1$ when $n=1$.
Thus the ordinary moments of every nonzero binomial
$X^a(A+BX^d)$ cannot both vanish at either adjacent pair of orders.
\end{corollary}
\begin{proof}
The preceding polynomial has smaller degree.
For the following polynomial, the
\href{BINOMIAL_IRREDUCIBILITY_REDUCTION.tex}{consecutive
nondivisibility theorem} excludes $P_{n,a,d}\mid P_{n+1,a,d}$.
An irreducible polynomial with neither divisibility possibility
is coprime to both neighbors over $\mathbb Q$; its Bezout identities
remain valid over $\mathbb C$.
The normalization
\[
 \mathcal L_1\bigl((X^a(A+BX^d))^m\bigr)
 = A^m(am)!P_{m,a,d}(B/A)
\]
gives the moment assertion when $AB\ne0$.
The monomial cases have nonzero positive-order moments.
\end{proof}

\section{Explicit ordinary quintic families}

\begin{corollary}\label{cor:quintic-residual-families}
For $a=2,d=3$, each row $(b,p)$ of the following table makes
$P_{bp^r,2,3}$ irreducible over $\mathbb Q$ for every $r\ge1$.
Both adjoining ordinary moment pairs are consequently coprime.
\[
 \begin{array}{c|c|c}
 b&p&\text{residual factor degrees}\\ \hline
 3&31&3\\
 3&37&3\\
 4&23&4\\
 4&29&1+3\\
 6&31&1+5
 \end{array}.
\]
\end{corollary}
\begin{proof}
The no-carry condition is satisfied because $p>5b$.
The literal monic residual factorizations are
\[
\begin{array}{c|c|l}
 b&p&\text{monic }P_{b,2,3}\pmod p\\ \hline
3&31&t^3+17t^2+13t+7\\
3&37&t^3+32t^2+22t+3\\
4&23&t^4+3t^3+12t^2+15t+21\\
4&29&(t+15)(t^3+13t^2+13t+3)\\
6&31&(t+6)(t^5+24t^4+22t^3+16t^2+27t+18)
\end{array}
\]
The accompanying exact verifier checks each product and verifies
irreducibility of each displayed factor $f$ of degree $h$ by
\[
 X^{p^h}\equiv X\pmod f,\qquad
 \gcd(f,X^{p^{h/q}}-X)=1
 \quad\text{for every prime }q\mid h.
 \tag{9}
\]
These conditions are a complete finite-field irreducibility criterion:
the first says that $f$ is squarefree and each irreducible factor
has degree dividing $h$; every proper divisor of $h$ divides $h/q$
for some prime $q\mid h$, so the gcd conditions exclude all proper
irreducible factors.
Theorem~\ref{thm:residual-index} proves the irreducible-residual rows,
and Corollary~\ref{cor:parity-residual} proves the two odd-factor rows.
Corollary~\ref{cor:residual-adjacent-pairs} gives the moment assertion.
\end{proof}

For example $n=4\cdot23^r$ yields still-unbounded subfamilies of the
previously remaining quintic progression: terminal $N=n$ when
$r\equiv1\pmod6$, and terminal $N=n+1$ when $r\equiv2\pmod6$.
These do not exhaust that progression.

\section{The quintic family: all-prime common-factor degrees}
Throughout this section specialize to $a=2,d=3,\beta=1$ and write
$P_n=P_{n,2,3}$. The no-carry hypothesis is imposed only where
explicitly stated; the all-prime degree criterion below does not need it.

Write
\[
 P_n(t)=\sum_{k=0}^n c_{n,k}t^k,\qquad
 c_{n,k}=\binom nk\frac{(2n+3k)!}{(2n)!}.
\]
All coefficients are positive integers and $c_{n,0}=1$.
A Newton slope below means a slope of the lower convex hull
of $(k,v_p(c_{n,k}))$, with adjacent collinear edges combined.
Its width is the horizontal length; a root associated with slope
$\lambda$ has valuation $-\lambda$.

\begin{lemma}[Elementary weighted product and root counts]\label{lem:quintic-weighted-newton-product}
For nonzero polynomials $f,g$ over an algebraic extension of
$\mathbb Q_p$, with $v_p(p)=1$, and rational $\nu$,
\[
 m_\nu(f):=\min_k\{v(a_k)+k\nu\},\qquad
 m_\nu(fg)=m_\nu(f)+m_\nu(g).
\]
Consequently Newton polygons add as slope multisets: the width at a
given slope is the number of roots, counted with multiplicity,
of the corresponding negative valuation.
For a polynomial with rational $p$-adic coefficients, the width
of every edge of reduced slope $a/b$ is divisible by $b$.
\end{lemma}
\begin{proof}
In a finite valued extension choose a common uniformizing power
which realizes the rational weight $\nu$ and the minimum.
Scale the variable and divide by a coefficient of minimum
valuation. The resulting polynomials have integral coefficients
and nonzero reductions; their product has nonzero reduction
because the residue polynomial ring is an integral domain.
This proves the weighted identity. Equivalently, one may use
the associated graded ring of the valuation.

Factor a polynomial into linear factors in a finite splitting
extension. For a linear factor $t-\alpha$ its weighted minimum is
$\min\{\nu,v(\alpha)\}$. Summing these piecewise-linear functions
shows that the slope break at $\nu=v(\alpha)$ has multiplicity
equal to the number of such roots. The same break computed
from the coefficient minimum is the corresponding horizontal
edge width. Finally, for rational $p$-adic coefficients every
Newton vertex has integer coordinates. An edge of width $w$
and reduced slope $a/b$ has integer height $aw/b$, so $b\mid w$.
\end{proof}

\begin{lemma}[Parity and full-degree exclusion]\label{lem:quintic-even-common-degree}
Every common rational factor of $P_n,P_{n+1}$ has even degree.
Also $P_n$ does not divide $P_{n+1}$; hence a nonconstant common
factor has degree $h$ with $2\le h<n$.
\end{lemma}
\begin{proof}
The even-degree assertion follows from
Lemma~\ref{lem:real-root-parity} at $a=2,d=3,\beta=1$.

Compare $c_n=[t]P_{n+1}-[t]P_n$ and
$L_n=[t^{n+1}]P_{n+1}/[t^n]P_n$. Direct cancellation gives
\[
 c_n=4(n+1)(2n+3)(4n+5),\qquad
 L_n=\frac{5\prod_{j=1}^4(5n+j)}{2(2n+1)}
\]
and the exact positive difference
\[
 L_n-c_n=
 \frac{3n(3n+2)(9n+7)(37n+17)}{2(2n+1)}>0.
\]
If $P_n$ divided $P_{n+1}$, their degrees and constant terms would
make the quotient $1+ct$. Its coefficient $c$ would equal
both $c_n$ and $L_n$, a contradiction.
\end{proof}

For each prime $p$, let $w_{n,p}(\lambda)$ be the width at slope
$\lambda$, or zero if that slope is absent.
Let $\Lambda_p(n)$ be the slopes occurring in both polygons.
Write each $\lambda\in\Lambda_p(n)$ in lowest terms as
$a_\lambda/b_\lambda$, $b_\lambda>0$, including $b_\lambda=1$
when $\lambda=0$. Define the finite set
\[
 {\mathcal D}_p(n)=
 \left\{\sum_{\lambda\in\Lambda_p(n)}b_\lambda j_\lambda:
 0\le j_\lambda\le
 \frac{\min\{w_{n,p}(\lambda),w_{n+1,p}(\lambda)\}}{b_\lambda},
 \quad j_\lambda\in\mathbb Z\right\}.                         \tag{D1}
\]
The upper bounds are integers by the weighted-product lemma.

\begin{theorem}[All-prime necessary common-factor degree criterion]\label{thm:quintic-all-prime-degree}
For every integer $n\ge1$, the degree of every nonconstant common
rational factor of $P_n,P_{n+1}$ belongs to
\[
 {\mathcal E}(n)=
 \{h:2\le h<n,\ h\text{ even}\}
 \cap\!\!\bigcap_{\substack{p\ {\rm prime}\\p\le5(n+1)}}
             {\mathcal D}_p(n).                              \tag{D2}
\]
In particular, $\mathcal E(n)=\varnothing$ implies
$\gcd(P_n,P_{n+1})=1$ over $\mathbb Q$ and over $\mathbb C$.
All inputs to (D2) are exact finite integer quantities:
\[
 v_p(c_{n,k})=
 v_p(n!)-v_p(k!)-v_p((n-k)!)
       +v_p((2n+3k)!)-v_p((2n)!).                            \tag{D3}
\]
This criterion has no no-carry hypothesis.
\end{theorem}
\begin{proof}
Let $G$ be a common rational factor of degree $h$.
At a prime $p$, all slopes of $G$ occur in both polygons,
by the weighted-product/root-count lemma.
At a common slope $a_\lambda/b_\lambda$, the number of roots
of $G$ is divisible by $b_\lambda$, and is at most each of the
two available widths. Summing these numbers gives exactly
the necessary condition $h\in\mathcal D_p(n)$.
The parity and full-degree lemma gives the first set in (D2).

For $p>5(n+1)$, every factorial and binomial coefficient defining
either polynomial is a $p$-adic unit. Both Newton polygons
are horizontal. Then $\mathcal D_p(n)=\{0,\ldots,n\}$,
so larger primes add no constraint.
Formula (D3) is immediate and its factorial valuations are
finite sums $\sum_{r\ge1}\lfloor m/p^r\rfloor$.

A nonconstant gcd over $\mathbb Q$ would have a degree in (D2).
If that set is empty, a rational Bezout identity exists and
remains a Bezout identity over $\mathbb C$.
\end{proof}

\paragraph{Missing all-index step.}
No proof is given that $\mathcal E(n)$ is empty for every $n$,
or for every $n\equiv1\pmod9$.
The sets in (D1) keep only necessary local slope/width data;
a surviving degree does not construct a rational common factor.
A finite or sparse collection of empty sets is not an
all-index result. The unresolved task is an unbounded
prime/digit argument or a stronger arithmetic invariant.

\begin{proposition}[An infinite family without any eligible no-carry part]\label{prop:quintic-no-carry-coverage-failure}
For every integer $s\ge0$, put
\[
 A=2^{3s+1},\qquad n=4A^4=2^{12s+6},\qquad N=n+1.
\]
Then $N\equiv2\pmod9$, but neither $n$ nor $N$ admits a
representation $m=bp^r$, $p$ prime and $r,b\ge1$, satisfying
the no-carry condition for lower exponent two and gap three:
\[
                    (2b\bmod p)+3b<p.                       \tag{D4}
\]
This also prevents applying any refinement which still requires
that same no-carry premise to either member of this family.
\end{proposition}
\begin{proof}
Condition (D4) implies $3b<p$, hence for $P=p^r$,
\[
                    P^2\ge pP>3bP=3m.                     \tag{D5}
\]
The order $n$ is a power of two and $p=2$ cannot satisfy (D4).

Sophie Germain's identity gives
\[
 N=(2A^2-2A+1)(2A^2+2A+1).
\]
The two factors are coprime: a common divisor is odd and
divides their difference $4A$, while each is congruent to one
modulo $A$. Thus every prime-power divisor $P$ of $N$ lies
entirely in one factor and is at most $2A^2+2A+1$.
For $A\ge2$,
\[
 (2A^2+2A+1)^2
 \le(13A^2/4)^2
 <12A^4<3N.
\]
This contradicts (D5). Finally $12s+6$ is divisible by six,
so $n\equiv1\pmod9$ and $N\equiv2\pmod9$.

This is failure of a sufficient criterion, not a counterexample
to consecutive coprimality. Other primes and Newton slopes
may still exclude the pair.
\end{proof}

\begin{corollary}[A complete two-prime certificate at $n=64$]\label{cor:quintic-n64-two-prime}
The pair $P_{64},P_{65}$ is coprime, although neither order has
an eligible no-carry prime-power part.
\end{corollary}
\begin{proof}
The preceding proposition with $s=0$ proves the no-carry
coverage failure. The complete lower-Newton slope/width lists are
\[
\begin{array}{c|c|l}
 p&\text{order}&(\text{slope},\text{width})\\\hline
 3&64&(1,1),(13/9,9),(40/27,27),(41/27,27)\\
 3&65&(1,1),(13/9,9),(40/27,27),(41/27,27),(4,1)\\
 5&64&(1/2,2),(3/5,5),(18/25,25),(19/25,25),(4/5,5),(1,2)\\
 5&65&(3/5,5),(18/25,25),(3/4,8),(19/25,25),(1,2).
\end{array}
\]
Each table is certified by the literal edge endpoints and all
coefficient supporting-line inequalities using (D3);
the exact checker also computes each coefficient
directly as an integer factorial quotient before valuation.
No degree greater than sixty-five or index prefix is used.

At $p=3$,
\[
 \mathcal D_3(64)=\{a+9b+27c:
        a,b\in\{0,1\},\ c\in\{0,1,2\}\}.
\]
Its positive even degrees below sixty-four are
$\{10,28,36,54\}$.
At $p=5$,
\[
 \mathcal D_5(64)=\{a+5b+25c:
       a\in\{0,1,2\},\ b\in\{0,1\},\ c\in\{0,1,2\}\}.
\]
None of the four remaining degrees belongs to this latter set.
The all-prime criterion proves the assertion, already using
only these two primes.
\end{proof}

\paragraph{Exact replay and scope.}
The exact checker controls factorial units, edge normalization,
factor-residual products, the complete small residual examples,
positive-Gamma root counts in bounded rational cases, and rejection
of violated carry or coefficient-valuation hypotheses.
For the single pair of orders $64,65$, it checks the four literal
Newton edge tables, every supporting-line inequality, and every
coefficient valuation independently from integer factorials and
Legendre's formula. It also checks the displayed exact algebra and
corruption rejection. Its literal Newton certificate is retained
separately from the replay receipt.
The proofs supply the unbounded residual, parity, all-prime
criterion and no-carry coverage-failure statements. No finite
index-prefix search establishes those quantifiers.
The ordinary quintic ray left by this reduction is closed by \href{QUINTIC_SCALAR_MINOR_SFC.tex}{the scalar-minor theorem}.
Internal replay is distinct from independent replay, external
review, and formal verification. No priority claim is made.
\end{document}
